\section{Interface Absorption Probability Using the Acoustic Mismatch Model}
\label{sec:interface_AMM}

A phonon that reaches the boundary between the sapphire substrate and a superconducting
film (Al, Nb or Ta) is either reflected back into the substrate or transmitted into the film,
where it can break Cooper pairs. The probability $P_a$ of this transmission enters the
simulation as the interface absorption probability. To estimate it from material properties
we use the loose-bond acoustic mismatch model of Every, Koos and
Wolfe~\cite{Every1984}, treating the substrate and the film as two isotropic half-spaces
joined by a layer of distributed springs. The complete derivation, including the coefficient
matrix of the boundary-value problem, is given in Appendix~\ref{app:interface_derivation};
here we state the model and discuss its results.

\paragraph{Elastic input.}
The model requires only the density $\rho$ and the isotropic Lam\'e parameters $\lambda$ and
$\mu$ of each material. These were computed in Sec.~\ref{TODO:sec_elastic} from the
second-order (and, for the anharmonic terms, third-order) elastic constants, and are listed in
Table~\ref{TODO:tab_lame}. They give the bulk wave speeds
\begin{equation}
  V_T=\sqrt{\mu/\rho},\qquad V_L=\sqrt{(\lambda+2\mu)/\rho},
\end{equation}
and the stress--strain relation
$\sigma_{ij}=\lambda\,\delta_{ij}\,\nabla\!\cdot\!\mathbf u+\mu(\partial_iu_j+\partial_ju_i)$
used in Sec.~\ref{TODO:sec_phonon_propagation} is the same one that enters the boundary
conditions here. Only the harmonic constants ($\lambda$, $\mu$) appear in the transmission
problem; the third-order coefficients ($\alpha$, $\beta$, $\gamma$) enter only the anharmonic
downconversion rate. Because $\lambda$ and $\mu$ are isotropic averages
(\todo{state the averaging scheme used for sapphire}), the sapphire substrate is treated
as isotropic in this chapter.

\paragraph{Boundary conditions.}
Let $z$ be the interface normal and let $u$ and $\sigma$ denote displacement and stress.
Across the interface the traction is continuous, and it is proportional to the jump in
displacement through the spring constants $K_\parallel$ (tangential) and $K_\perp$ (normal),
\begin{equation}
  \sigma_{xz}=K_\parallel\,\Delta u_x,\qquad \sigma_{zz}=K_\perp\,\Delta u_z .
\end{equation}
The coupling is expressed through the dimensionless parameters
\begin{equation}
  \alpha_\parallel=\frac{K_\parallel}{\rho_cV_{T,c}\,\omega},\qquad
  \alpha_\perp=\frac{K_\perp}{\rho_cV_{L,c}\,\omega},
  \label{eq:alpha_omega}
\end{equation}
where $\rho_c$, $V_{T,c}$ and $V_{L,c}$ refer to the substrate and $\omega$ is the phonon
angular frequency. If the stiffness $K$ is frequency independent, $\alpha\propto1/\omega$,
so low-energy phonons see a stronger coupling than high-energy ones. The limit
$\alpha\to\infty$ is a welded interface, which reproduces the classical acoustic mismatch
result, and $\alpha\to0$ is a free surface.

\paragraph{Reflection and transmission of a single wave.}
An incident L or SV wave in the substrate generates four outgoing waves: reflected L and SV
in the substrate and transmitted L and SV in the film. Their angles follow from Snell's law,
$\sin\theta_j/V_j=\sin\theta_i/V_i$, and their amplitudes follow from solving the four
boundary conditions (Appendix~\ref{app:interface_derivation}). The fraction of the incident
power carried by each outgoing wave is
\begin{equation}
  W_j=|\Gamma_j|^2\,\frac{\rho_jV_j\cos\theta_j}{\rho_{\rm inc}V_{\rm inc}\cos\theta_{\rm inc}},
  \label{eq:power_coeff}
\end{equation}
where $\Gamma_j$ is the displacement amplitude of wave $j$ relative to the incident wave.
Evanescent waves carry no power. The four fractions sum to one at every angle, as expected
for an elastic interface that dissipates no energy. For normal incidence the transmitted
power has the closed form
\begin{equation}
  W_{\rm trans}=\frac{4r}{(1+r)^2+r^2/\alpha^2},\qquad r=\frac{Z_m}{Z_c},
  \label{eq:normal_spring_main}
\end{equation}
with $Z=\rho V$ the acoustic impedance of the relevant wave type; it reduces to
$4r/(1+r)^2$ for a welded interface (Eq.~\eqref{eq:normal_spring}).

Figure~\ref{sec:fig_interface_angular} illustrates the angular dependence for sapphire/Ta.
For an SV wave incident from sapphire, the reflected L wave becomes evanescent beyond the
critical angle
\begin{equation}
  \theta_c=\arcsin\frac{V_{T,c}}{V_{L,c}}=36.4^\circ ,
\end{equation}
and the pseudo-surface wave at $\theta_c$ produces a cusp in the transmitted power that is
visible for all couplings. A weaker coupling lowers the transmission at all angles; for
$\alpha_\parallel=\alpha_\perp=0.3$ the transmission at normal incidence drops from about
$0.98$ (welded) to about $0.2$.

\begin{figure}[H]
\centering
\includegraphics[width=0.95\textwidth]{Figures/interface_angular.pdf}
\caption{Sapphire/Ta interface, isotropic elastic constants.
(a)~Reflection and refraction of an incident L wave at the interface. Each incident wave
produces four outgoing waves, L and SV in sapphire and L and SV in Ta, with angles fixed by
Snell's law, $\sin\theta_j/v_j=\sin\theta_i/v_i$. The inset sketches the loose-bond (spring)
boundary condition, in which the stresses are proportional to the displacement jump,
$\sigma_{zz}=K_\perp\Delta u_z$ and $\sigma_{xz}=K_\parallel\Delta u_x$.
(b)~Inset of (c): partition of the power of an incident L wave among the four outgoing waves
as a function of the incidence angle $\theta_i$, for $\alpha_\parallel=1.0$ and
$\alpha_\perp=1.1$; the four contributions sum to one at every angle.
(c)~Squared amplitude $|\Gamma_{T\to L}|^2$ of the reflected L wave for a unit incident SV
wave. For a free surface, the pseudo-surface wave gives a sharp peak at the critical angle
$\theta_c=\arcsin(v_T/v_L)\approx36^\circ$; this peak is progressively suppressed as the
coupling to the metal increases.
(d)~Power transmitted into Ta for an incident SV wave, for a welded interface and for the two
loosely bonded cases indicated in the legend. The cusp at $\theta_c$ is the same resonance
seen in~(c).}
\label{sec:fig_interface_angular}
\end{figure}

\paragraph{Angle- and mode-averaged probability.}
Phonons arrive from the substrate with an approximately isotropic distribution, so the
relevant quantity is the flux-weighted angular average of $W_{\rm trans}$ for each incident
mode, with weight $\cos\theta\sin\theta$, combined over the L, SV and SH modes with the Debye
weights $V^{-3}$ (one L and two transverse modes). The result,
$\langle W\rangle_{\rm Debye}$ (Appendix~\ref{app:interface_derivation}), is listed in
Table~\ref{tab:interface_P} for Cu, Al, Nb and Ta on sapphire. For $\alpha\gtrsim2$ the
probability saturates at the welded value, which depends only on the acoustic impedance
mismatch; for $\alpha\lesssim1$ it falls steeply with decreasing $\alpha$
(Fig.~\ref{sec:fig_interface_vs_alpha}).

\begin{table}[H]
\centering
\caption{Debye- and flux-averaged transmission probability $\langle W\rangle_{\rm Debye}$
from sapphire into the metal film. The ``realistic'' column uses
$\alpha_\parallel=1.0$, $\alpha_\perp=1.1$ (Cu-fitted values of
Ref.~\cite{Every1984}), ``order unity'' uses $\alpha_\parallel=\alpha_\perp=1.0$ and
``very loose'' uses $\alpha_\parallel=\alpha_\perp=0.3$. \todo{The coupling parameters are
not fitted for Al, Nb and Ta; the three columns are a sensitivity range, not a prediction.}}
\label{tab:interface_P}
\begin{tabular}{lcccc}
\toprule
Film & welded & realistic & order unity & very loose\\
\midrule
Cu & 0.937 & 0.871 & 0.865 & 0.401\\
Al & 0.829 & 0.809 & 0.807 & 0.566\\
Nb & 0.937 & 0.876 & 0.871 & 0.412\\
Ta & 0.856 & 0.760 & 0.754 & 0.277\\
\bottomrule
\end{tabular}
\end{table}

\begin{figure}[H]
\centering
\includegraphics[width=0.8\textwidth]{Figures/interface_vs_alpha.pdf}
\caption{Debye- and flux-averaged transmission probability from sapphire into Cu, Al and Ta
as a function of the coupling $\alpha=\alpha_\parallel=\alpha_\perp$. Dotted lines mark the
welded limit of each metal. The upper axis converts $\alpha$ to phonon energy under the
assumptions that $\alpha=1$ at $2.48$~meV ($\approx600$~GHz) and that the bond stiffness $K$
is frequency independent; both are assumptions of this work and not measured values.}
\label{sec:fig_interface_vs_alpha}
\end{figure}

\paragraph{Limitations.}
Three assumptions limit how literally these numbers should be read.
First, the coupling parameters $\alpha_\parallel$ and $\alpha_\perp$ are not determined by the
elastic constants; the values used are those fitted for Cu by Every \emph{et
al.}~\cite{Every1984} or are scanned (\todo{measure or cite $K$ for Al/Nb/Ta on sapphire}).
Second, SH waves are given the tangential coupling, $\alpha_{SH}=\alpha_\parallel$, because
their motion is purely tangential; this controls the SH contribution to the Ta penalty.
Third, sapphire is treated as isotropic with averaged elastic constants, although it is
anisotropic, and the phonon-energy dependence of $\alpha$ through $\omega$ is an assumption.
The densities used for the isotropic velocities are \todo{confirm: 3986, 2700, 8570 and
16650~kg\,m$^{-3}$ for sapphire, Al, Nb and Ta}.
